\documentclass{article}
\usepackage[utf8]{inputenc}
\usepackage{amsmath}
\usepackage{geometry}
\geometry{margin=2cm}
\begin{document}

\section*{Análise da Questão 161 - ENEM 2024 - Matemática}

\subsection*{Enunciado resumido}
Um tanque prismático retangular possui dois anteparos verticais fixados na sua base, dividindo sua base em três retângulos congruentes. Os anteparos possuem alturas diferentes, \(\frac{H}{4}\) e \(\frac{H}{2}\), e são instalados em paredes opostas, ortogonais entre si. O tanque é preenchido com água por um duto no teto com vazão constante, levando 12 horas para encher completamente.

\subsection*{Objetivo}
Identificar o gráfico que mostra a maior altura de coluna de água no tanque em função do tempo durante o enchimento.

\subsection*{Solução técnica}
\begin{itemize}
  \item O tanque tem altura total \(H\), largura total \(L\), com a base dividida em 3 partes iguais de largura \(\frac{L}{3}\).
  \item Os anteparos limitam a altura máxima da água em duas secções: uma a \(\frac{H}{4}\) e outra a \(\frac{H}{2}\).
  \item Devido à vazão constante, o volume de água aumenta linearmente no tempo.
  \item Inicialmente, a água enche o compartimento com anteparo menor até \(\frac{H}{4}\), maior altura inicial da coluna de água.
  \item Depois, a água preenche o segundo compartimento até \(\frac{H}{2}\), enquanto o primeiro mantém altura constante.
  \item Finalmente, a água enche o terceiro compartimento até altura \(H\).
  \item O gráfico da maior altura da coluna de água apresenta três segmentos:
    \begin{enumerate}
      \item Crescimento linear até \(\frac{H}{4}\).
      \item Platô constante entre \(\frac{H}{4}\) e \(\frac{H}{2}\).
      \item Crescimento linear final até \(H\).
    \end{enumerate}
  \item O tempo total, 12 horas, distribui-se proporcionalmente aos volumes das etapas.
\end{itemize}

\subsection*{Explicação didática}
Este problema estimula a compreensão das relações entre volume, área e altura de líquido, além de interpretar o efeito de barreiras verticais dentro do tanque.

\subsubsection*{Passos para resolução}
\begin{enumerate}
  \item Identificar a divisão da base do tanque em três partes iguais.
  \item Reconhecer as alturas dos anteparos \(\frac{H}{4}\) e \(\frac{H}{2}\).
  \item Notar que o tanque está sendo enchido com vazão constante e, portanto, volume aumenta linearmente no tempo.
  \item Entender que cada etapa do enchimento limita a altura máxima da coluna de água em função dos anteparos.
  \item Mapear as fases do enchimento no gráfico da altura máxima em função do tempo.
\end{enumerate}

\subsubsection*{Erros comuns}
\begin{itemize}
  \item Misturar a ordem dos compartimentos que se enchem.
  \item Supor enchimento simultâneo nos três compartimentos.
  \item Não interpretar corretamente os platôs no gráfico.
  \item Ignorar a vazão constante e o aumento linear do volume.
\end{itemize}

\subsubsection*{Recomendações pedagógicas}
\begin{itemize}
  \item Usar modelos físicos ou software para simular o enchimento.
  \item Trabalhar as relações entre volume, área e altura em sólidos geométricos.
  \item Explorar leitura e interpretação de gráficos em contextos físicos.
\end{itemize}

\subsection*{Análise dos distractores}
\begin{itemize}
  \item \textbf{Alternativa A}: Erro no tempo de duração do platô, embora o formato geral pareça correto.
  \item \textbf{Alternativa C}: Platô demasiado longo com subida rápida muito curta, contrariando o volume proporcional.
  \item \textbf{Alternativa D}: Vários platôs sem subida linear adequada, contradizendo a vazão constante.
  \item \textbf{Alternativa E}: Apresenta quedas na altura, o que não é possível no enchimento.
\end{itemize}

\subsection*{Perfil da questão}
Esta questão exige interpretação gráfica, raciocínio geométrico, noções de proporção e modelagem do processo de enchimento com vazão constante, sendo adequada para o nível médio do ENEM.

\end{document}


\end{document}